% ── Aspectos Administrativos ──────────────────────────────────────────────────
\section{Aspectos Administrativos}

\subsection{Cronograma}

El proyecto se organiza en \textbf{cuatro sprints} de dos semanas cada uno, más una
fase inicial de planificación (\textit{Sprint 0}), alineados con las dos entregas
formales del semestre, tal como se describe en la sección de metodología. El trabajo
inició el \textbf{martes 19 de mayo} con el Sprint 0, el cual comprende el análisis
del TUPA vigente, el levantamiento de requerimientos, la elaboración del plan de
proyecto y el prototipado dinámico, culminando con el \textbf{primer entregable}
el \textbf{28 de mayo de 2026}.

El equipo realiza reuniones de coordinación diarias (\textit{Daily Stand-up}) de
10 minutos para sincronizar avances, identificar impedimentos y planificar el
trabajo del día, complementadas con las ceremonias formales de Scrum (\textit{Sprint
Planning}, \textit{Sprint Review} y \textit{Sprint Retrospective}) al inicio y
cierre de cada sprint.

\begin{figure}[H]
\centering
\caption{Diagrama de Gantt --- Semestre 2026-I}
\vspace{0.4cm}
\begin{ganttchart}[
  hgrid,
  vgrid={*{1}{draw=gray!25}},
  x unit=0.43cm,
  y unit title=0.55cm,
  y unit chart=0.50cm,
  title height=1,
  title label font=\scriptsize\bfseries,
  bar label font=\tiny\raggedleft,
  bar label node/.append style={align=right},
  group label font=\small\bfseries\color{unsaacRed},
  milestone label font=\tiny\itshape,
  bar/.style={fill=unsaacRed!65,rounded corners=1.5pt},
  bar incomplete/.style={fill=unsaacRed!25,rounded corners=1.5pt},
  group/.style={draw=unsaacRed,fill=unsaacRed!15},
  group peaks height=0.15,
  milestone/.style={fill=unsaacRed,shape=diamond,inner sep=2pt},
  today=10,
  today rule/.style={draw=unsaacRed!80,dashed,line width=0.7pt},
  today label={\scriptsize Hoy},
  today label node/.append style={below=2pt,font=\tiny,color=unsaacRed},
]{1}{22}
%--- Encabezados ---
\gantttitle{Semestre 2026-I (22 semanas)}{22} \\
\gantttitle{Mayo}{4}\gantttitle{Junio}{5}\gantttitle{Julio}{8}\gantttitle{Agosto}{5} \\
\gantttitlelist{1,...,22}{1} \\

%--- Sprint 0 ---
\ganttgroup{Sprint 0 --- Planificación y Requerimientos}{1}{2} \\
\ganttbar{Análisis del TUPA vigente}{1}{1} \\
\ganttbar{Levantamiento de requerimientos}{1}{2} \\
\ganttbar{Redacción del plan de proyecto}{1}{2} \\
\ganttbar{Prototipado dinámico (Figma)}{2}{2} \\
\ganttmilestone{Entrega 1 (28/05)}{2} \\[grid]

%--- Sprint 1 ---
\ganttgroup{Sprint 1 --- Módulo de Consulta TUPA}{3}{6} \\
\ganttbar{Implementación de consulta de procedimientos}{3}{4} \\
\ganttbar{Autenticación y gestión de usuarios}{3}{4} \\
\ganttbar{Layout base y diseño responsivo}{3}{5} \\
\ganttbar{Configuración de BD y migraciones}{3}{3} \\
\ganttbar{Pruebas unitarias del módulo}{4}{6} \\[grid]

%--- Sprint 2 ---
\ganttgroup{Sprint 2 --- Registro y Seguimiento}{7}{10} \\
\ganttbar{Formulario de registro de trámite}{7}{8} \\
\ganttbar{Seguimiento de expedientes}{7}{9} \\
\ganttbar{Sistema de notificaciones}{8}{10} \\
\ganttbar{Validación de documentos}{8}{9} \\[grid]

%--- Sprint 3 ---
\ganttgroup{Sprint 3 --- Panel Administrativo}{11}{14} \\
\ganttbar{Panel de control administrativo}{11}{13} \\
\ganttbar{Módulo de reportes y estadísticas}{12}{14} \\
\ganttbar{Mejoras de accesibilidad}{13}{14} \\[grid]

%--- Sprint 4 ---
\ganttgroup{Sprint 4 --- Integración y Cierre}{15}{21} \\
\ganttbar{Integración de módulos}{15}{17} \\
\ganttbar{Corrección de defectos}{15}{18} \\
\ganttbar{Pruebas de regresión}{16}{19} \\
\ganttbar{Documentación técnica}{17}{20} \\
\ganttbar{Despliegue en la nube}{19}{21} \\
\ganttmilestone{Entrega 2 (1.${}^{\circ}$ sem. agosto)}{22}
\end{ganttchart}
\end{figure}

\subsection{Presupuesto}

El presupuesto estimado considera los recursos necesarios para el desarrollo del
sistema web durante el semestre 2026-I con un equipo de cuatro integrantes a lo
largo de los cuatro sprints del proyecto. Los costos se agrupan en cinco categorías:
recursos humanos (valorizados a una tarifa referencial de mercado), equipamiento
informático, software y licencias, servicios de hosting y dominio, y materiales de
oficina e imprevistos. La estimación sigue criterios conservadores, priorizando
herramientas con planes gratuitos o de bajo costo, dado que se trata de un proyecto
académico autofinanciado por los integrantes del grupo.

\begin{table}[H]
\centering
\small
\caption{Presupuesto estimado del proyecto (en Soles peruanos --- S/.)}
\begin{tabular}{C{0.5cm} L{5.2cm} C{1.8cm} C{1.6cm} R{2.0cm}}
\toprule
\rowcolor{unsaacRed}
\color{white}\textbf{N.º} &
\color{white}\textbf{Concepto} &
\color{white}\textbf{Unidad} &
\color{white}\textbf{Cantidad} &
\color{white}\textbf{Costo (S/.)} \\
\midrule
\multicolumn{5}{l}{\textbf{A. Recursos Humanos}} \\
\midrule
1 & Programador iniciante (4 integrantes)
    & Horas & 384 & \num{7680,00} \\
\multicolumn{4}{r}{\textit{Subtotal A}} & \textbf{\num{7680,00}} \\
\midrule
\multicolumn{5}{l}{\textbf{B. Equipamiento}} \\
\midrule
2 & Computadoras portátiles (costo de uso semestral, 4 unid.)
    & Global & 4 & \num{800,00} \\
3 & Servidor de desarrollo (VPS, 4 meses)
    & Mes & 4 & \num{120,00} \\
4 & Periféricos (mouse, teclado, hub USB)
    & Global & 1 & \num{80,00} \\
\multicolumn{4}{r}{\textit{Subtotal B}} & \textbf{\num{1000,00}} \\
\midrule
\multicolumn{5}{l}{\textbf{C. Software y Licencias}} \\
\midrule
5 & Claude Pro (plan más barato, 4 meses)
    & Mes & 4 & \num{300,00} \\
6 & Cuenta GitHub (plan gratuito)
    & --- & --- & \num{0,00} \\
7 & Cuenta Figma (plan gratuito)
    & --- & --- & \num{0,00} \\
\multicolumn{4}{r}{\textit{Subtotal C}} & \textbf{\num{300,00}} \\
\midrule
\multicolumn{5}{l}{\textbf{D. Hosting y Dominio}} \\
\midrule
8 & Registro de dominio (.com / .pe, 1 año)
    & Año & 1 & \num{50,00} \\
9 & Hosting en la nube (Railway / Render --- gratuito)
    & --- & --- & \num{0,00} \\
\multicolumn{4}{r}{\textit{Subtotal D}} & \textbf{\num{50,00}} \\
\midrule
\multicolumn{5}{l}{\textbf{E. Materiales de Oficina e Imprevistos}} \\
\midrule
10 & Materiales de oficina (lapiceros, folders, etc.)
    & Global & 1 & \num{30,00} \\
11 & Imprevistos (5\,\% del subtotal de B + C + D)
    & --- & --- & \num{67,50} \\
\multicolumn{4}{r}{\textit{Subtotal E}} & \textbf{\num{97,50}} \\
\midrule
\rowcolor{lightGray}
\multicolumn{4}{r}{\textbf{TOTAL GENERAL}} & \textbf{S/.\,\num{9127,50}} \\
\bottomrule
\end{tabular}
\end{table}

\begin{supuesto}
\small\textbf{Nota:} Los costos de recursos humanos (A) se valorizan a una tarifa
referencial de S/.\,20,00 por hora para programadores iniciantes, considerando
2 horas diarias, 3 días por semana (martes, jueves y sábado), durante las
16 semanas del semestre (384 horas totales). El equipamiento (B) considera el
costo de uso y desgaste semestral de las laptops personales de los integrantes.
El software (C) incluye una suscripción compartida de Claude Pro. El total
asciende a \textbf{S/.\,\num{9127,50}}, financiado íntegramente por los
integrantes del grupo.
\end{supuesto}
